\setcounter{section}{14}

  \zwischenueberschrift{Bentuk diferensial pada manifold}

\inputdefinition
{}
{

Misalkan $M$ suatu
\definitionsverweis {manifold diferensiabel}{}{.}
Suatu
\definitionswortpraemath {k}{ bentuk diferensial }{}
\zusatzklammer {atau $k$-\definitionswort {bentuk}{} atau \definitionswort {bentuk berderajat}{} $k$} {} {}
adalah suatu
\definitionsverweis {penampang}{}{}
dalam
\definitionsverweis {produk eksterior}{}{}
ke-$k$ dari
\definitionsverweis {bundel kotangen}{}{,}
yaitu suatu pemetaan
\maabbeledisp {\omega} { M } { \bigwedge^k T^*M
} { P } { \omega(P)
} {,}
dengan

\mavergleichskette
{\vergleichskette
{ \omega(P)
}
{ \in }{ \bigwedge^k T_P^*M
}
{  }{
}
{    }{
}
{   }{
}
}
{}{}{.}

}

Kita melambangkan himpunan semua $k$-bentuk pada $M$ dengan

\mathdisp {{ \mathcal E }^{ k } ( M  )} { . }

\inputbemerkung
{}
{

Jadi, suatu $k$-\definitionsverweis {bentuk}{}{}
memasangkan kepada setiap titik $P$ pada manifold sebuah elemen dari
\mathl{\bigwedge^k T^*_PM}{}. Berdasarkan
Korolari 57.9 (Aljabar Linear (Osnabrück 2024-2025))
dan Teorema 58.4 (Aljabar Linear (Osnabrück 2024-2025)),
hal ini sama dengan suatu
\definitionsverweis {pemetaan multilinear}{}{}
\definitionsverweis {selang-seling}{}{}
\maabbdisp {} {T_PM \times \cdots \times T_PM } { \R
} {.}
Pemetaan yang bersesuaian ini juga kita lambangkan dengan
\mathl{\omega(P)}{;} untuk $k$
\definitionsverweis {vektor singgung}{}{}

\mavergleichskette
{\vergleichskette
{ v_1 , \ldots , v_k
}
{ \in }{ T_PM
}
{  }{
}
{    }{
}
{   }{
}
}
{}{}{}
dengan demikian,

\mathdisp {\omega(P)(v_1 , \ldots , v_k)} {  }

adalah suatu bilangan riil. Di sini muncul dua jenis argumen yang sama sekali berbeda: di satu pihak titik pada manifold, dan di pihak lain elemen-elemen ruang singgung di titik tersebut. Ketergantungan pada vektor-vektor singgung relatif sederhana karena pemetaannya multilinear selang-seling, sedangkan ketergantungan pada titik manifold dapat serumit apa pun. Karena produk-produk eksterior bundel kotangen sendiri merupakan manifold menurut
Soal 14.2,
kita dapat langsung membicarakan bentuk diferensial yang kontinu atau
\zusatzklammer {jika $M$ merupakan manifold $C^2$} {} {}
diferensiabel.

Untuk

\mavergleichskette
{\vergleichskette
{ k
}
{ = }{ 0
}
{  }{
}
{    }{
}
{   }{
}
}
{}{}{}
ruang kotangen hanya muncul secara formal; suatu $0$-bentuk tidak lain adalah sebuah fungsi
\maabb {f} { M } {\R
} {.}
Suatu $1$-bentuk
\zusatzklammer {juga disebut
\definitionswortenp{bentuk Pfaff}{}} {} {}
memasangkan suatu bilangan riil kepada setiap titik $P$ dan setiap vektor singgung pada titik tersebut. Untuk

\mavergleichskette
{\vergleichskette
{k
}
{ > }{ n
}
{ = }{ \operatorname{dim} \, M
}
{    }{
}
{   }{
}
}
{}{}{}
produk eksterior ke-$k$ adalah ruang nol sehingga sama sekali tidak ada bentuk taktrivial berderajat ini. Kasus yang sangat penting adalah

\mavergleichskette
{\vergleichskette
{ k
}
{ = }{ n
}
{ = }{ \operatorname{dim} \, M
}
{    }{
}
{   }{
}
}
{}{}{.}
Dalam hal ini, produk eksterior ke-$n$ memiliki rank $1$
\zusatzklammer {artinya, dimensinya $1$ di setiap titik} {} {}
dan suatu penampang di dalamnya secara lokal dideskripsikan oleh satu fungsi saja. Gambaran yang berguna ialah bahwa, untuk $n$ vektor singgung, bilangan
\mathl{\omega(P)(v_1 , \ldots , v_n)}{} menyatakan volume
\zusatzklammer {\anfuehrung{berorientasi}{}} {} {}
dari paralelotop yang direntang oleh vektor-vektor tersebut di ruang singgung. Gambaran ini juga membantu untuk $k$ yang lebih kecil: dengan
\mathl{\omega(P) (v_1 , \ldots , v_k)}{} kita dapat menghitung volume berdimensi $k$ dari paralelotop yang dibangkitkan oleh $k$ vektor singgung. Gambaran ini akan dibuat tepat ketika kita mengintegralkan pada suatu
\definitionsverweis {submanifold tertutup}{}{}
berdimensi $k$.

}



\inputfaktbeweis
{Mannigfaltigkeit/Differentialform/Elementare Eigenschaften/Fakt}
{Lemma}
{}
{

\faktsituation {Misalkan $M$ suatu
\definitionsverweis {manifold diferensiabel}{}{}
dan, untuk

\mavergleichskette
{\vergleichskette
{ k
}
{ \in }{ \N
}
{  }{
}
{    }{
}
{   }{
}
}
{}{}{},
misalkan
\mathl{{ \mathcal E }^{ k } (  M   )}{} himpunan semua
$k$-\definitionsverweis {bentuk}{}{}
pada $M$.}
\faktuebergang {Maka berlaku sifat-sifat berikut.}
\faktfolgerung {\aufzaehlungfuenf{Dengan operasi-operasi alaminya,
\mathl{{ \mathcal E }^{ k } (  M   )}{} membentuk
\definitionsverweis {ruang vektor riil}{}{.}
}{Untuk suatu bentuk diferensial

\mavergleichskette
{\vergleichskette
{ \omega
}
{ \in }{
{ \mathcal E }^{ k } (  M   )
}
{  }{
}
{    }{
}
{   }{
}
}
{}{}{}
dan suatu fungsi
\maabbdisp {f} { M } {\R
} {},

\mavergleichskette
{\vergleichskette
{ f \omega
}
{ \in }{
{ \mathcal E }^{ k } (  M   )
}
{  }{
}
{    }{
}
{   }{
}
}
{}{}{,}
dengan $f \omega$ didefinisikan oleh

\mavergleichskettedisp
{\vergleichskette
{ (f \omega)(P)
}
{ \defeq} { f(P) \omega (P)
}
{ } {
}
{ } {
}
{ } {
}
}
{}{}{}.
}{Setiap
\definitionsverweis {fungsi diferensiabel}{}{}
$C^1$
\maabbdisp {f} { M } {\R
} {}
melalui
\definitionsverweis {pemetaan singgung}{}{}

\mathl{Tf}{} mendefinisikan suatu
$1$-\definitionsverweis {bentuk diferensial}{}{}
\maabbeledisp {df} { M } { T^*M
} { P } { T_Pf
} {,}
dengan ruang singgung dari $\R$ di
\mathl{f(P)}{} diidentifikasi dengan $\R$. Hal ini menghasilkan suatu pemetaan
\maabbeledisp {d} {C^1(M,\R)} {
{ \mathcal E }^{  1 } (  M   )
} { f } {df
} {.}
}{Jika

\mavergleichskette
{\vergleichskette
{ M
}
{ \subseteq }{ \R^m
}
{  }{
}
{    }{
}
{   }{
}
}
{}{}{}
merupakan himpunan bagian terbuka, maka di bawah identifikasi

\mavergleichskettedisp
{\vergleichskette
{ TM
}
{ \cong} { M \times \R^m
}
{ } {
}
{ } {
}
{ } {
}
}
{}{}{}
pemetaan pada (3) sama dengan
\maabbeledisp {} { M } { M \times ( \R^m)^*
} { P } { (P, \left(  D f   \right)_{ P } \left(   -  \right) )
} {.}
}{Pemetaan $d$ pada (3) adalah
$\R$-\definitionsverweis {linear}{}{.}
}}
\faktzusatz {}
\faktzusatz {}

}
{

\teilbeweis {}{}{}
{(1) dan (2) langsung mengikuti dari peninjauan titik demi titik.}
{} \teilbeweis {}{}{}
{(3). Untuk setiap titik

\mavergleichskette
{\vergleichskette
{ P
}
{ \in }{ M
}
{  }{
}
{    }{
}
{   }{
}
}
{}{}{},

\maabbdisp {T_Pf} { T_PM } { T_{f(P)} \R \cong \R
} {}
merupakan suatu
\definitionsverweis {pemetaan linear}{}{}
menurut Lemma 9.10  (3),
dan dengan demikian merupakan elemen dari
\mathl{T_P^* M}{,} yang kita lambangkan dengan
\mathl{(df)_P}{}. Oleh karena itu, pemasangan
\mathl{P \mapsto (df)_P}{} merupakan suatu
\definitionsverweis {bentuk diferensial}{}{.}}
{}
\teilbeweis {}{}{}
{(4) mengikuti dari
Lemma 9.10  (1).}
{}
\teilbeweis {}{}{}
{(5). Pemetaan pada (3) ditentukan untuk setiap titik

\mavergleichskette
{\vergleichskette
{ P
}
{ \in }{ M
}
{  }{
}
{    }{
}
{   }{
}
}
{}{}{}
pada setiap lingkungan terbuka. Karena itu, kita boleh menganggap bahwa

\mavergleichskette
{\vergleichskette
{ M
}
{ \subseteq }{ \R^m
}
{  }{
}
{    }{
}
{   }{
}
}
{}{}{}
merupakan himpunan terbuka, sehingga pernyataannya mengikuti dari (4) dan
Proposisi 45.6 (Analisis (Osnabrück 2021-2023)).}
{}

}

Untuk suatu himpunan terbuka

\mavergleichskette
{\vergleichskette
{ V
}
{ \subseteq }{ \R^n
}
{  }{
}
{    }{
}
{   }{
}
}
{}{}{},
tersedia fungsi-fungsi koordinat
\maabbdisp {x_j} {V} {\R
} {},
yang, jika diberikan suatu peta, dapat dipindahkan ke suatu himpunan bagian terbuka dari manifold. Di setiap titik

\mavergleichskette
{\vergleichskette
{ Q
}
{ \in }{ V
}
{  }{
}
{    }{
}
{   }{
}
}
{}{}{},
bentuk-bentuk

\mathbed {dx_j} {}
 {j=1 , \ldots , n} {}
 {} {} {} {,}
membentuk basis ruang kotangen di $Q$. Basis ini tidak lain adalah
\definitionsverweis {basis dual}{}{}
dari basis standar di ruang ambien $\R^n$, yang digunakan sebagai ruang singgung di seluruh $V$. Untuk suatu himpunan bagian beranggotakan $k$

\mavergleichskettedisp
{\vergleichskette
{ J
}
{ =} { \{j_1 , \ldots , j_k\}
}
{ \subseteq} { \{1 , \ldots , n\}
}
{ } {
}
{ } {
}
}
{}{}{},
kita tetapkan

\mavergleichskettedisp
{\vergleichskette
{  dx_J
}
{ =} {  dx_{j_1} \wedge \ldots \wedge dx_{j_k}
}
{ } {
}
{ } {
}
{ } {
}
}
{}{}{,}
yang merupakan $k$-bentuk yang sangat sederhana pada $V$. Untuk setiap titik

\mavergleichskette
{\vergleichskette
{ Q
}
{ \in }{ V
}
{  }{
}
{    }{
}
{   }{
}
}
{}{}{},
berlaku

\mavergleichskettedisp
{\vergleichskette
{ dx_J(Q)
}
{ =} { (dx_{j_1} \wedge \ldots \wedge dx_{j_k} )(Q)
}
{ =} { dx_{j_1}(Q) \wedge \ldots \wedge dx_{j_k} (Q)
}
{ =} { e_{j_1}^* \wedge \ldots \wedge e_{j_k}^*
}
{ } {
}
}
{}{}{.}
Aksi bentuk ini pada

\mavergleichskette
{\vergleichskette
{  v_1 \wedge \ldots \wedge v_k
}
{ \in }{ \bigwedge^k T_QV
}
{  }{
}
{    }{
}
{   }{
}
}
{}{}{}
, menurut
Teorema 58.4 (Aljabar Linear (Osnabrück 2024-2025))
, diberikan oleh

\mavergleichskettedisp
{\vergleichskette
{ \left(  e_{j_1}^* \wedge \ldots \wedge e_{j_k}^*  \right)  (v_1 \wedge \ldots \wedge v_k)
}
{ =} { \det  \left(  e_{j_i}^*(v_\ell)  \right)_{i \ell}
}
{ =} { \det  \left(  (v_\ell)_{j_i}  \right)_{i \ell}
}
{ } {
}
{ } {
}
}
{}{}{.}
Menurut
Teorema 58.1 (Aljabar Linear (Osnabrück 2024-2025)),
nilai-nilai bentuk diferensial
\zusatzklammer {dengan \mathlk{j_1 <   \ldots < j_k}{}} {} {}

\mavergleichskettedisp
{\vergleichskette
{ dx_J
}
{ =} { dx_{j_1} \wedge \ldots \wedge dx_{j_k}
}
{ } {
}
{ } {
}
{ } {
}
}
{}{}{}
membentuk suatu
\definitionsverweis {basis}{}{}
dari
\mathl{\bigwedge^k T_Q^*V}{}
untuk setiap titik $Q$. Oleh karena itu, setiap
\definitionsverweis {bentuk diferensial}{}{} pada $V$ yang berderajat $k$

\mavergleichskette
{\vergleichskette
{ \omega
}
{ \in }{
{ \mathcal E }^{ k } ( V  )
}
{  }{
}
{    }{
}
{   }{
}
}
{}{}{}
dapat ditulis secara tunggal sebagai

\mavergleichskettedisp
{\vergleichskette
{ \omega
}
{ =} { \sum_{J,\, { \# \left(  J \right) } = k} f_J dx_J
}
{ } {
}
{ } {
}
{ } {
}
}
{}{}{}
dengan fungsi-fungsi yang ditentukan secara tunggal
\maabbdisp {f_J} {V} {\R
} {.}



\inputfaktbeweis
{Mannigfaltigkeit/Differentialform/Lokale Beschreibung/Fakt}
{Lemma}
{}
{

\faktsituation {Misalkan $M$ suatu
\definitionsverweis {manifold diferensiabel}{}{}
dan

\mavergleichskette
{\vergleichskette
{ U
}
{ \subseteq }{ M
}
{  }{
}
{    }{
}
{   }{
}
}
{}{}{}
suatu
\definitionsverweis {himpunan bagian terbuka}{}{}
dengan peta
\maabbdisp {\alpha} { U } { V
} {},
dan misalkan

\mavergleichskette
{\vergleichskette
{ V
}
{ \subseteq }{ \R^n
}
{  }{
}
{    }{
}
{   }{
}
}
{}{}{}
terbuka. Misalkan
\maabbdisp {x_j} { U } {\R
} {}
fungsi-fungsi koordinat yang bersesuaian,

\mavergleichskette
{\vergleichskette
{ 1
}
{ \leq }{ j
}
{ \leq }{ n
}
{    }{
}
{   }{
}
}
{}{}{.}}
\faktfolgerung {Maka setiap
\definitionsverweis {bentuk diferensial}{}{} pada $U$ yang berderajat $k$

\mavergleichskette
{\vergleichskette
{ \omega
}
{ \in }{
{ \mathcal E }^{ k } (  U   )
}
{  }{
}
{    }{
}
{   }{
}
}
{}{}{}
dapat ditulis secara tunggal sebagai

\mavergleichskettedisp
{\vergleichskette
{ \omega
}
{ =} { \sum_{J,\, { \# \left(   J  \right) } = k}  f_J dx_J
}
{ } {
}
{ } {
}
{ } {
}
}
{}{}{}
dengan fungsi-fungsi yang ditentukan secara tunggal
\maabbdisp {f_J} { U } {\R
} {.}}
\faktzusatz {}
\faktzusatz {}

}
{

Hal ini mengikuti dari
Teorema 58.1 (Aljabar Linear (Osnabrück 2024-2025)).

}



\inputfaktbeweis
{Mannigfaltigkeit/Pfaffsche Differentialform zu Funktion/Beschreibung mit partiellen Ableitungen/Fakt}
{Korollar}
{}
{

\faktsituation {Misalkan $M$ suatu
\definitionsverweis {manifold diferensiabel}{}{}
dan

\mavergleichskette
{\vergleichskette
{U
}
{ \subseteq }{ M
}
{  }{
}
{    }{
}
{   }{
}
}
{}{}{}
suatu
\definitionsverweis {himpunan bagian terbuka}{}{}
dengan peta
\maabbdisp {\alpha} { U } { V
} {},
dan misalkan

\mavergleichskette
{\vergleichskette
{V
}
{ \subseteq }{\R^n
}
{  }{
}
{    }{
}
{   }{
}
}
{}{}{}
terbuka. Misalkan
\maabbdisp {x_j} { U } {\R
} {}
fungsi-fungsi koordinat yang bersesuaian,

\mavergleichskette
{\vergleichskette
{ 1
}
{ \leq }{j
}
{ \leq }{n
}
{    }{
}
{   }{
}
}
{}{}{.}
Misalkan
\maabbdisp {f} { U } {\R
} {}
suatu
\definitionsverweis {fungsi diferensiabel}{}{.}}
\faktfolgerung {Maka $1$-\definitionsverweis {bentuk diferensial}{}{}
$df$ yang bersesuaian mempunyai representasi\zusatzfussnote {Turunan-turunan
\mathl{{ \frac{   \partial  f  }{  \partial  x_j  } }}{} diperkenalkan dalam kuliah kedelapan} {.} {}

\mavergleichskettedisp
{\vergleichskette
{ df
}
{ =} { \sum_{j = 1}^n { \frac{   \partial   f   }{  \partial   x_j   } } dx_j
}
{ } {
}
{ } {
}
{ } {
}
}
{}{}{.}}
\faktzusatz {}
\faktzusatz {}

}
{

Kita dapat langsung menganggap bahwa semuanya berlangsung pada himpunan terbuka

\mavergleichskette
{\vergleichskette
{V
}
{ \subseteq }{\R^n
}
{  }{
}
{    }{
}
{   }{
}
}
{}{}{}.
Untuk setiap titik

\mavergleichskette
{\vergleichskette
{Q
}
{ \in }{ V
}
{  }{
}
{    }{
}
{   }{
}
}
{}{}{},
berlaku kesamaan berikut antara
\definitionsverweis {bentuk-bentuk linear}{}{}
pada $\R^n$,

\mavergleichskettealign
{\vergleichskettealign
{(df)_Q
}
{ =} { \left(D f \right)_{ Q }
}
{ =} { \left(  { \frac{   \partial   f   }{  \partial   x_1   } } ( Q)   , \,   \ldots  , \,   { \frac{   \partial   f   }{  \partial   x_n   } } ( Q)                 \right)
}
{ =} { { \frac{   \partial   f   }{  \partial   x_1   } } ( Q) dx_1 + \cdots + { \frac{   \partial   f   }{  \partial   x_n   } } ( Q) dx_n
}
{ } {
}
}
{}
{}{.}

}

  \zwischenueberschrift{Penarikan balik bentuk diferensial}

\inputdefinition
{}
{

Misalkan
\mathkor {} {L} {dan} {M} {}
\definitionsverweis {dua manifold diferensiabel}{}{}
dan misalkan
\maabbdisp {\varphi} { L } { M
} {}
suatu
\definitionsverweis {pemetaan diferensiabel}{}{.}
Misalkan $\omega$ suatu $k$-\definitionsverweis {bentuk diferensial}{}{}
pada $M$. Maka pemetaan yang menentukan suatu $k$-bentuk pada $L$ diberikan oleh

\mathdisp {(P ,v_1 , \ldots , v_k) \longmapsto   \omega(\varphi(P), T_P(\varphi)(v_1) , \ldots ,  T_P(\varphi)(v_k) )} {  }

Pemetaan ini merupakan
\definitionsverweis {pemetaan selang-seling}{}{}.
Bentuk yang bersesuaian dengannya disebut, terhadap $\varphi$, \definitionswort {bentuk tarik balik}{} berderajat $k$. Bentuk ini dilambangkan dengan

\mathdisp {\varphi^* \omega} {  }.

}



\inputfaktbeweis
{Mannigfaltigkeit/Differenzierbare Abbildung/Zurückziehen von Differentialformen/Elementare Eigenschaften/Fakt}
{Lemma}
{}
{

\faktsituation {Misalkan
\mathkor {} {L} {dan} {M} {}
\definitionsverweis {dua manifold diferensiabel}{}{}
dan
\maabbdisp {\varphi} { L } { M
} {}
suatu
\definitionsverweis {pemetaan diferensiabel}{}{.}}
\faktuebergang {Maka
\definitionsverweis {penarikan balik bentuk diferensial}{}{}
memenuhi sifat-sifat berikut.}
\faktfolgerung {\aufzaehlungvier{Untuk suatu fungsi

\mavergleichskette
{\vergleichskette
{ f
}
{ \in }{ \operatorname{Map} (M,\R)
}
{ = }{
{ \mathcal E }^{  0 } (  M   )
}
{    }{
}
{   }{
}
}
{}{}{},
berlaku

\mavergleichskette
{\vergleichskette
{ \varphi^*(f)
}
{ = }{ f \circ \varphi
}
{  }{
}
{    }{
}
{   }{
}
}
{}{}{.}
}{Pemetaan-pemetaan
\maabbeledisp {\varphi^*} {
{ \mathcal E }^{ k } (  M   ) } {
{ \mathcal E }^{ k } (  L   )
} { \omega} { \varphi^* \omega
} {,}
adalah
$\R$-\definitionsverweis {linear}{}{.}
}{Jika

\mavergleichskette
{\vergleichskette
{ L
}
{ \subseteq }{ M
}
{  }{
}
{    }{
}
{   }{
}
}
{}{}{}
merupakan suatu
\definitionsverweis {submanifold terbuka}{}{,}
dan pemetaan di atas adalah inklusi, maka penarikan balik suatu bentuk diferensial $\omega$ tidak lain adalah pembatasan
\mathl{\omega \vert_L}{} pada himpunan bagian ini.
}{Misalkan $N$ suatu manifold diferensiabel lain dan
\maabbdisp {\psi} { M } { N
} {}
suatu pemetaan diferensiabel lain. Maka berlaku

\mavergleichskettedisp
{\vergleichskette
{ ( \psi \circ \varphi )^*(\omega)
}
{ =} { \varphi^*(\psi^*(\omega))
}
{ } {
}
{ } {
}
{ } {
}
}
{}{}{}
untuk setiap bentuk diferensial

\mavergleichskette
{\vergleichskette
{ \omega
}
{ \in }{
{ \mathcal E }^{ k } (  N   )
}
{  }{
}
{    }{
}
{   }{
}
}
{}{}{.}
}}
\faktzusatz {}
\faktzusatz {}

}
{

\teilbeweis {}{}{}
{(1) langsung mengikuti dari
Definisi 14.6.}
{}
\teilbeweis {}{}{}
{(2). Untuk bentuk-bentuk diferensial
\mathkor {} {\omega_1} {dan} {\omega_2} {}
dan skalar-skalar

\mavergleichskette
{\vergleichskette
{ a,b
}
{ \in }{\R
}
{  }{
}
{    }{
}
{   }{
}
}
{}{}{},
kita harus menunjukkan bahwa

\mavergleichskette
{\vergleichskette
{ \varphi^*(a \omega_1+ b \omega_2)
}
{ = }{ a \varphi^* \omega_1 + b \varphi^* \omega_2
}
{  }{
}
{    }{
}
{   }{
}
}
{}{}{}.
Kesamaan bentuk diferensial semacam ini berarti bahwa kesamaan tersebut berlaku di setiap titik

\mavergleichskette
{\vergleichskette
{ P
}
{ \in }{ L
}
{  }{
}
{    }{
}
{   }{
}
}
{}{}{}
dan untuk setiap $k$-tupel vektor singgung

\mavergleichskette
{\vergleichskette
{ v_1 , \ldots , v_k
}
{ \in }{ T_PL
}
{  }{
}
{    }{
}
{   }{
}
}
{}{}{}.
Karena itu, pernyataannya mengikuti dari

\mavergleichskettealignhandlinks
{\vergleichskettealignhandlinks
{ \left(  \varphi^*(a \omega_1 + b \omega_2)  \right) \left(  P, v_1 , \ldots , v_k  \right)
}
{ =} { (a \omega_1 + b \omega_2) \left(  \varphi(P), T_P(\varphi)(v_1) , \ldots ,  T_P(\varphi)(v_k)  \right)
}
{ =} { a \omega_1 \left(  \varphi(P), T_P(\varphi)(v_1) , \ldots ,  T_P(\varphi)(v_k)  \right)
}
{ \,\,\,\, \,} {+ b \omega_2 \left(  \varphi(P), T_P(\varphi)(v_1) , \ldots ,  T_P(\varphi)(v_k)  \right)
}
{ =} { \left(  a  \varphi^* \omega_1  \right)(P, v_1 , \ldots , v_k)  + \left(  b \varphi^* \omega_2  \right) (P, v_1 , \ldots , v_k)
}
}
{}
{}{.}}
{}
\teilbeweis {}{}{}
{(3) langsung mengikuti dari definisi.}
{}
\teilbeweis {}{}{}
{(4). Misalkan

\mavergleichskette
{\vergleichskette
{ Q
}
{ \in }{ L
}
{  }{
}
{    }{
}
{   }{
}
}
{}{}{,}

\mavergleichskette
{\vergleichskette
{ u_1 , \ldots , u_k
}
{ \in }{ T_QL
}
{  }{
}
{    }{
}
{   }{
}
}
{}{}{},
dan $\omega$ suatu $k$-bentuk pada $N$. Maka, dengan menggunakan
Lemma 9.10  (4),
berlaku

\mavergleichskettealigndrucklinks
{\vergleichskettealigndrucklinks
{(( \psi \circ \varphi)^* (\omega))(Q, u_1 , \ldots , u_k)
}
{ =} { \omega \left(  ( \psi \circ \varphi) (Q), T_Q( \psi \circ \varphi)(u_1) , \ldots , T_Q(  \psi \circ \varphi) (u_k)  \right)
}
{ =} { \omega \left(  \psi ( \varphi (Q)), \left(  T_{\varphi(Q)}\psi  \right)  ((T_Q\varphi)(u_1)) , \ldots , \left(  T_{\varphi(Q)} \psi  \right) ((T_Q\varphi)(u_k))  \right)
}
{ =} { \left(  \psi^*(\omega)  \right) \left(  \varphi(Q),T_Q(\varphi)(u_1) , \ldots , T_Q(\varphi)(u_k)  \right)
}
{ =} { \left(  \varphi^* \left(  \psi^*(\omega)  \right)   \right) (Q,u_1 , \ldots , u_k)
}
}
{}{}{,}
dan inilah pernyataannya.}
{}

}



\inputfaktbeweis
{Offene Mengen/Differenzierbare Abbildung/Zurückziehen von Differentialformen/In Koordinaten/Fakt}
{Lemma}
{}
{

\faktsituation {Misalkan
\mathkor {} {U \subseteq \R^n} {dan} {V \subseteq \R^m} {}
\definitionsverweis {himpunan-himpunan bagian terbuka}{}{,}
dengan koordinat masing-masing dilambangkan
\mathl{x_1 , \ldots , x_n}{} dan
\mathl{y_1 , \ldots , y_m}{}. Misalkan
\maabbdisp {\varphi} { U } { V
} {}
suatu
\definitionsverweis {pemetaan diferensiabel}{}{}
dan misalkan $\omega$ suatu
$k$-\definitionsverweis {bentuk diferensial}{}{}
pada $V$ dengan representasi

\mavergleichskettedisp
{\vergleichskette
{ \omega
}
{ =} { \sum_{I,\, { \# \left(   I  \right) } = k}  f_I  dy_I
}
{ } {
}
{ } {
}
{ } {
}
}
{}{}{,}
dengan
\maabb {f_I} { V } {\R
} {}
merupakan
\definitionsverweis {fungsi-fungsi}{}{.}}
\faktfolgerung {Maka
\definitionsverweis {bentuk tarik balik}{}{}
mempunyai representasi

\mavergleichskettealignhandlinks
{\vergleichskettealignhandlinks
{ \varphi^* \omega
}
{ =} { \sum_{I,\, { \# \left(   I  \right) } = k}  \left(  f_I \circ \varphi  \right)  \left( \sum_{J,\, { \# \left(   J  \right) } = k} \det  \left(  \left(  { \frac{   \partial  \varphi_i   }{  \partial   x_j   } }  \right)_{ i \in I, j \in J}  \right)  dx_J \right)
}
{ =} { \sum_{J,\, { \# \left(   J  \right) } = k}   \left( \sum_{I,\, { \# \left(   I  \right) } = k} \left(  f_I \circ \varphi  \right)  \det  \left(  \left(  { \frac{   \partial  \varphi_i   }{  \partial   x_j   } }  \right)_{ i \in I, j \in J}  \right)  \right)  dx_J
}
{ } {
}
{ } {
}
}
{}{}{.}}
\faktzusatz {}
\faktzusatz {}

}
{

Kesamaan kedua didasarkan pada pengaturan ulang yang sederhana. Berdasarkan
Lemma 14.7 (2),
kita dapat membatasi diri pada kasus

\mavergleichskette
{\vergleichskette
{ \omega
}
{ = }{ f_Idy_I
}
{  }{
}
{    }{
}
{   }{
}
}
{}{}{}.
Kita tetapkan

\mavergleichskette
{\vergleichskette
{ f
}
{ = }{ f_I
}
{  }{
}
{    }{
}
{   }{
}
}
{}{}{}
dan boleh menganggap

\mavergleichskette
{\vergleichskette
{ I
}
{ = }{ \{ 1 , \ldots , k \}
}
{  }{
}
{    }{
}
{   }{
}
}
{}{}{}.
Kita menunjukkan kesamaan antara kedua $k$-bentuk pada $U$ dengan menunjukkan bahwa keduanya menghasilkan nilai yang sama untuk setiap titik

\mavergleichskette
{\vergleichskette
{ P
}
{ \in }{ U
}
{  }{
}
{    }{
}
{   }{
}
}
{}{}{}
dan setiap produk eksterior
\mathl{e_{j_1 } \wedge \ldots \wedge e_{j_k }}{} dengan

\mavergleichskette
{\vergleichskette
{ 1
}
{ \leq }{ j_1
}
{ < }{ \ldots
}
{ <   }{ j_k
}
{ \leq  }{ n
}
}
{}{}{}.
Di satu pihak,

\mavergleichskettealigndrucklinks
{\vergleichskettealigndrucklinks
{ \varphi^*( f dy_1 \wedge \ldots \wedge dy_k)(P,e_{j_1 } \wedge \ldots \wedge e_{j_k }  )
}
{ =} { ( f dy_1 \wedge \ldots \wedge dy_k)(\varphi(P),T_P(\varphi)(e_{j_1 }) \wedge \ldots \wedge T_P(\varphi)(e_{j_k }))
}
{ =} { f(\varphi(P)) (dy_1 \wedge \ldots \wedge dy_k)( \begin{pmatrix}  { \frac{   \partial  \varphi_1   }{  \partial   x_{j_1 }  } } (P )  \\\vdots\\  { \frac{   \partial  \varphi_m   }{  \partial   x_{j_1 }  } } (P )   \end{pmatrix} \wedge \ldots \wedge \begin{pmatrix}  { \frac{   \partial  \varphi_1   }{  \partial   x_{j_k }  } } (P )  \\\vdots\\  { \frac{   \partial  \varphi_m   }{  \partial   x_{j_k }  } } (P )   \end{pmatrix} )
}
{ =} { f(\varphi(P)) \cdot \det  \left(  (dy_i) \begin{pmatrix}  { \frac{   \partial  \varphi_1   }{  \partial   x_{j_\ell}  } } (P )  \\\vdots\\  { \frac{   \partial  \varphi_m   }{  \partial   x_{j_\ell}  } } (P )   \end{pmatrix}  \right)_{1 \leq i, \ell \leq k}
}
{ =} { f(\varphi(P)) \cdot \det  \left(  { \frac{   \partial  \varphi_i   }{  \partial   x_{j_\ell}  } } (P )  \right)_{1 \leq i, \ell \leq k}
}
}
{}{}{.}
Di pihak lain, jika jumlah tersebut diterapkan pada
\mathl{e_{j_1 } \wedge \ldots \wedge e_{j_k }}{},
maka

\mavergleichskette
{\vergleichskette
{ dx_J (e_{j_1 } \wedge \ldots \wedge e_{j_k })
}
{ = }{ 0
}
{  }{
}
{    }{
}
{   }{
}
}
{}{}{}
kecuali apabila

\mavergleichskette
{\vergleichskette
{ J
}
{ = }{ \{j_1 , \ldots , j_k\}
}
{  }{
}
{    }{
}
{   }{
}
}
{}{}{,}
dan dalam kasus itu nilainya $1$, sehingga diperoleh nilai yang sama.

}



\inputfaktbeweis
{Offene Mengen/Differenzierbare Abbildung/Zurückziehen von Volumenform auf gleichdimensionalem Raum/In Koordinaten/Fakt}
{Korollar}
{}
{

\faktsituation {Misalkan

\mavergleichskette
{\vergleichskette
{ U,V
}
{ \subseteq }{ \R^n
}
{  }{
}
{    }{
}
{   }{
}
}
{}{}{}
\definitionsverweis {himpunan-himpunan bagian terbuka}{}{,}
dengan koordinat masing-masing dilambangkan
\mathl{x_1 , \ldots , x_n}{} dan
\mathl{y_1 , \ldots , y_n}{}. Misalkan
\maabbdisp {\varphi} { U } { V
} {}
suatu
\definitionsverweis {pemetaan diferensiabel}{}{}
dan misalkan $\omega$ suatu
$n$-\definitionsverweis {bentuk diferensial}{}{}
pada $V$ dengan representasi

\mavergleichskettedisp
{\vergleichskette
{ \omega
}
{ =} { f dy_1 \wedge \ldots \wedge dy_n
}
{ } {
}
{ } {
}
{ } {
}
}
{}{}{.}}
\faktfolgerung {Maka
\definitionsverweis {bentuk tarik balik}{}{}
mempunyai representasi

\mavergleichskettedisp
{\vergleichskette
{ \varphi^* \omega
}
{ =} { (f \circ \varphi) \cdot \det  \left(  \left(  { \frac{   \partial  \varphi_i   }{  \partial   x_j   } }  \right)_{1 \leq  i, j \leq n}  \right) dx_1 \wedge \ldots \wedge dx_n
}
{ } {
}
{ } {
}
{ } {}
}
{}{}{.}}
\faktzusatz {}
\faktzusatz {}

}
{

Hal ini langsung mengikuti dari
Lemma 14.8.

}
Dengan demikian, fungsi-fungsi yang mendeskripsikan suatu bentuk diferensial mempunyai perilaku transformasi yang sama dengan kerapatan berorientasi pada suatu peta. Untuk kerapatan ukuran positif, faktor determinannya harus diganti dengan nilai mutlak determinan.



\inputfaktbeweis
{Differentialform/Lokal/Zurückziehen unter partiell konstanter Abbildung/Fakt}
{Korollar}
{}
{

\faktsituation {Misalkan

\mavergleichskette
{\vergleichskette
{ U
}
{ \subseteq }{ \R^n
}
{  }{
}
{    }{
}
{   }{
}
}
{}{}{}
dan

\mavergleichskette
{\vergleichskette
{ V
}
{ \subseteq }{ \R^m
}
{  }{
}
{    }{
}
{   }{
}
}
{}{}{}
\definitionsverweis {himpunan-himpunan bagian terbuka}{}{,}
dengan koordinat masing-masing dilambangkan
\mathl{x_1 , \ldots , x_n}{} dan
\mathl{y_1 , \ldots , y_m}{}. Misalkan
\maabbdisp {\varphi} { U } { V
} {}
suatu
\definitionsverweis {pemetaan diferensiabel}{}{}
dengan
\mathl{\varphi_{i_0 }}{} konstan untuk suatu

\mavergleichskette
{\vergleichskette
{ i_0
}
{ \in }{ \{1 , \ldots , m\}
}
{  }{
}
{    }{
}
{   }{
}
}
{}{}{},
dan misalkan $\omega$ suatu
$k$-\definitionsverweis {bentuk diferensial}{}{}
pada $V$ dengan representasi

\mavergleichskettedisp
{\vergleichskette
{ \omega
}
{ =} { f dy_I
}
{ } {
}
{ } {
}
{ } {
}
}
{}{}{}
dengan

\mavergleichskette
{\vergleichskette
{ i_0
}
{ \in }{ I
}
{  }{
}
{    }{
}
{   }{
}
}
{}{}{.}}
\faktfolgerung {Maka

\mavergleichskette
{\vergleichskette
{ \varphi^*\omega
}
{ = }{ 0
}
{  }{
}
{    }{
}
{   }{
}
}
{}{}{.}}
\faktzusatz {}
\faktzusatz {}

}
{

Menurut
Lemma 14.8,
berlaku

\mavergleichskettedisp
{\vergleichskette
{ \varphi^* \omega
}
{ =} { \sum_{J,\, { \# \left(   J  \right) } = k} (f \circ \varphi) \cdot \det  \left(  \left(  { \frac{   \partial  \varphi_i   }{  \partial   x_j   } }  \right) _{ i \in I, j \in J}  \right) dx_J
}
{ } {
}
{ } {
}
{ } {}
}
{}{}{.}
Karena

\mavergleichskettedisp
{\vergleichskette
{ { \frac{   \partial  \varphi_{i_0 }  }{  \partial   x_j   } }
}
{ =} { 0
}
{ } {
}
{ } {
}
{ } {
}
}
{}{}{}
untuk semua

\mavergleichskette
{\vergleichskette
{ j
}
{ \in }{ J
}
{  }{
}
{    }{
}
{   }{
}
}
{}{}{,}
untuk setiap $J$, matriks tersebut mempunyai satu baris bernilai $0$, sehingga semua determinannya bernilai $0$.

}
